\section{Excellence}\label{sec:excellence}

\HECriterion{Excellence --- aspects to be taken into account}{%
\begin{itemize}
  \item Clarity and pertinence of the project's objectives, and the extent to
  which the proposed work is ambitious and goes beyond the state of the art.
  \item Soundness of the proposed methodology, including the underlying
  concepts, models, assumptions, interdisciplinary approaches, appropriate
  consideration of the gender dimension in research and innovation content,
  and the quality of open science practices, including sharing and management
  of research outputs and engagement of citizens, civil society and end users
  where appropriate.
\end{itemize}
These aspects are considered only to the extent that the proposed work is within
the scope of the work programme topic.}

\subsection{Objectives and ambition\HEBudget{e.g. 4 pages}}
\label{sec:objectives}
\begin{heguidance}
\begin{itemize}
  \item Briefly describe the objectives of the proposed work. Explain their
  pertinence to the work programme topic, how they can be measured and verified,
  and why they are realistically achievable.
  \item Describe how the project goes beyond the state of the art and the
  extent to which the proposed work is ambitious. Identify any exceptional
  ground-breaking R\&I, novel concepts and approaches, new products, services,
  or business and organisational models. Where relevant, compare with products
  and services already available on the market and refer to patent or publication
  searches carried out.
  \item Position the proposed work in terms of R\&I maturity: from idea to
  application, or from lab to market. Where applicable, state the Technology
  Readiness Levels at the start and end of the project.
\end{itemize}
Interpret advances beyond the state of the art in the light of the project's
positioning: expectations for lower-TRL RIAs differ from those for high-TRL IAs.
The suggested section lengths are indicative; the whole-document limit applies.
\end{heguidance}

\HEParagraph{Objectives}
\HEPlaceholder{State the objectives, their connection to the topic, measurable
success indicators, and means of verification.}

\HEParagraph{Ambition and advance beyond the state of the art}
\HEPlaceholder{Describe the research gap, relevant baselines, the proposed
advance, and the project's starting and target maturity. Cite supporting work.}
\HEPlaceholder{Citation example: \cite{SumersYNG2024}. Replace this example
with references relevant to the proposal.}

\subsection{Methodology\HEBudget{e.g. 14 pages}}\label{sec:methodology}

\HEParagraph{Overall methodology}
\begin{heguidance}
Describe and explain the overall methodology, including the concepts, models and
assumptions underpinning the work. Explain how it enables the project to deliver
its objectives, identify important methodological challenges, and explain how
they will be overcome. Present a narrative here; detailed tasks and work packages
belong in section 3.1. Suggested allocation: about 10 pages.
\end{heguidance}
\HEPlaceholder{Explain the methodological choices, assumptions, validation
strategy, evidence to be collected, and links to each objective.}

\HEParagraph{Links to other research and innovation activities}
\begin{heguidance}
Describe national or international research and innovation activities whose
results will feed into the project and explain how those links will be
established. Suggested allocation: about 1 page.
\end{heguidance}
\HEPlaceholder{Identify the relevant activities and explain the concrete
mechanism by which their results will contribute to this project.}

\HEParagraph{Interdisciplinarity}
\begin{heguidance}
Explain how expertise and methods from different disciplines will be brought
together and integrated in pursuit of the objectives. If an interdisciplinary
approach is unnecessary, justify that assessment. Suggested allocation: about
half a page.
\end{heguidance}
\HEPlaceholder{Explain the integration of disciplines, including interfaces
between their methods and contributions.}

\HEParagraph{Social sciences and humanities}
\begin{heguidance}
For topics whose work programme indicates a need to integrate social sciences
and humanities, show the role of these disciplines in the project or justify why
they are not relevant to the proposed work. Suggested allocation: about half a
page.
\end{heguidance}
\HEPlaceholder{Explain the contribution of social sciences and humanities
where required by the topic, or provide the relevant justification.}

\HEParagraph{Gender dimension in research and innovation content}
\begin{heguidance}
Describe how sex and/or gender analysis is taken into account in the research
and innovation content. If it is not relevant, justify that assessment. This
section is mandatory except for topics explicitly identified in the work
programme as not requiring integration of the gender dimension. Suggested
allocation: about 1 page.

This concerns the content of the research and innovation activities, rather
than gender balance in the project team. Sex analysis concerns biological
characteristics; gender analysis concerns social and cultural factors.
\end{heguidance}
\HEPlaceholder{Explain where sex or gender may affect the research questions,
methods, data, analysis, validation or outcomes, and how this will be addressed.}

\HEParagraph{Open science and research output management}
\begin{heguidance}
Describe how appropriate open science practices form an integral part of the
methodology. Explain how the practices and their implementation suit the work
and increase its chances of achieving the objectives. If none are appropriate,
provide a justification. Suggested allocation: about 1 page.

Open science involves open cooperative work and systematic sharing of knowledge
and tools as early and widely as possible. Relevant practices include early
and open sharing through preregistration, registered reports, preprints or
crowdsourcing; research output management; reproducibility measures; open access
to publications, data, software, models, algorithms and workflows; open peer
review; and co-creation with citizens, civil society and end users where
appropriate. Outreach, communication, dissemination and exploitation belong
in section 2.2.

A funded project must deliver a detailed data management plan (DMP) by month 6
and revise it towards the end of the project. Explain how generated or collected
data and other research outputs will be findable, accessible, interoperable and
reusable (FAIR). The Horizon Europe Programme Guide provides further guidance
on open science and research data management.
\end{heguidance}
\HEPlaceholder{Specify practices, responsibilities, repositories, access
conditions, reproducibility measures, engagement activities where appropriate,
and the approach to FAIR research output management.}

